\section{Open problems and lessons for the Thorin programme}\label{ch:open}

Section~\ref{ch:path} located the gap on the path to RH: positivity uniform in the height, first as the GGC property at a single basepoint below one (step~7 of Table~\ref{tab:ch25:path}) and then across the strip. The programme of computations below is aimed at that gap, and the lessons that follow it are drawn for anyone working on the Thorin route to RH, including the first author's own programme \cite{P26}. Three of the five work packages of the programme, the twisted Riesz filters, the positivity margins against random matrix theory and the almost GGC defect, are taken up in Sections~\ref{ch:riesz} and \ref{ch:almost}; the open problems on the reciprocal zeta function are stated in Section~\ref{rz:sec:open}, those on the basepoint one and the critical strip in Section~\ref{ag:sec:exactly}, and those of Pakes's framework in Section~\ref{pk:sec:open}, and we do not repeat them.

\subsection{What computation can and cannot do on the path}\label{sec:open}\label{st:sec:role}

On the path of Section~\ref{sec:ch25:path}, for $\alpha\ge1$ the Thorin density is positive, the L\'evy form of $\log\xi(\alpha+s)$ converges, the normal scale mixtures are extended GGCs and the tilted P\'olya kernels are complete in the sense of Wiener, all with no hypothesis (Sections~\ref{ch:basepoint} to \ref{ch:bpeta}), and for $\frac12<\alpha<1$ each of these, read as the GGC property of the shifted clock $T_\alpha$ and as convergence, the extended GGC property and completeness at every basepoint to the right of $\alpha$, is equivalent to a zero free half plane $\re s>\alpha$.

Computation cannot supply uniformity in height (Section~\ref{sec:ch25:contrib}): every finite computation leaves room for zeros above the computed range, and a single off line zero, however high, stops $T_\alpha$ from being a GGC at every basepoint $\alpha$ between $\frac12$ and its real part, whether or not it makes the Thorin density negative near its ordinate (Section~\ref{sec:ch13:visibility}). What computation can do falls under four aims. It can extend and certify: prove positivity of the Thorin density on explicit windows, at every resolution, with interval arithmetic. It can test intermediate conjectures: the four questions of Section~\ref{ag:sec:exactly} (a lower bound for the background density of order $\log\gamma$ outside a sparse set, thin strip density bounds, an explicit bound for the defect of the canonical factorisation and positivity at one basepoint) all have numerically testable consequences. It can calibrate: measure how much room positivity has, and compare with random matrix predictions, so that a proof knows which inequalities are sharp. And it can cross check: compute the same quantities from the primes, through the criteria of Riesz and Salem and the compensated L\'evy form, independently of evaluations of $\zeta$.

Five work packages serve these aims. The first, twisted Riesz functions as an arithmetic filter bank, is the cross check of Section~\ref{ri:sec:filters}; the second, certified Thorin densities at height, is the subject of Section~\ref{st:sec:wp2}; the third and fourth, the positivity margins against the circular unitary ensemble and explicit tables of the almost GGC defect, are described in Sections~\ref{ag:sec:margins} and \ref{ag:sec:factor}; and the fifth, testing stronger conjectures before pursuing them, is the subject of Section~\ref{st:sec:wp5}. Table~\ref{st:tab:mile} lists their milestones and the resources each needs. The computations that remain from the chapters of the book fall under the same aims, and Sections~\ref{st:sec:wp2} to \ref{st:sec:wp5} take them together with the packages.

\begin{table}[htbp]
\centering\small
\begin{tabular}{p{0.1\linewidth}p{0.47\linewidth}p{0.33\linewidth}}
\toprule
package & milestone & resources\\
\midrule
1 & $R(x)$ and $R_h(x)$ for $x\le10^{16}$, $h$ on a grid to $500$, compared with the explicit formula & compiled sieve, days of CPU\\
2 & certified $\rho_\varepsilon$ at $\varepsilon=10^{-1},10^{-2},10^{-3}$ on windows at heights $10^6$ to $10^{10}$ & Arb, weeks of CPU\\
3 & margin distributions from Odlyzko's tables near the $10^{12}$th \cite{Odlyzko87}, $10^{21}$st and $10^{22}$nd zeros \cite{Odlyzko01,OdlyzkoTables}; CUE extreme value fit & hours\\
4 & explicit almost GGC defect tables & hours\\
5 & numerical tests of the four questions of Section~\ref{ag:sec:exactly} and of HCM below one & ongoing\\
\bottomrule
\end{tabular}
\caption{Work packages and milestones of the computational programme. Package 1 is described in Section~\ref{ri:sec:filters}, where $R$ is Riesz's function and $R_h$ its twist; packages 3 and 4 are described in Sections~\ref{ag:sec:margins} and \ref{ag:sec:factor}.}
\label{st:tab:mile}
\end{table}

\subsection{Certified Thorin densities at height}\label{st:sec:wp2}

At resolution $\varepsilon=\alpha-\frac12$ the Thorin density $\rho_\varepsilon(\theta)=\frac1\pi\re\frac{\xi'}{\xi}(\alpha+i\theta)$ is a smooth function of $\theta$. Its positivity on a window of ordinates is what the absence of zeros with $\beta>\alpha$ predicts there; since an off line zero need not make the density negative (Section~\ref{sec:ch13:visibility}), a certificate for the window should also count such zeros. Two computations do this. The first evaluates $\rho_\varepsilon$ on a grid with rigorous error bounds, using the ball arithmetic Riemann Siegel and Euler Maclaurin evaluators of Arb \cite{Arb}, exposed in python flint, at moderate height and the P\'olya route where it is faster, and bounds the variation between grid points by the derivative. The second counts the zeros to the right of the line $\re s=\alpha$ through the zero density identity $M_\varepsilon(T)+\tau_\varepsilon(T)=N(T)-2\nu(\frac12+\varepsilon,T)$ of Proposition~\ref{prop:zerodensity}, which returns an integer from a smooth computation at any resolution. It is the count of Section~\ref{ch:turing}, which at $\varepsilon=\frac12$ is Backlund's method. Section~\ref{sec:ch9:count} followed the argument in steps of $\varepsilon/8$ ($\varepsilon/4$ at $\varepsilon=0.005$) without a bound on its change between grid points, so its counts are computed, not certified; with such a bound the count would certify a whole window, at an estimated cost of a few hundred evaluations per unit height at resolution $\varepsilon\approx0.1$. The deliverables are certified positivity of $\rho_\varepsilon$ for $\varepsilon\in\{10^{-1},10^{-2},10^{-3}\}$ on windows of length $10^4$ at heights $10^6$, $10^8$ and $10^{10}$, together with the minimum density on each window, which feeds the margin statistics of Section~\ref{ag:sec:margins}. Below the verified height $3\times10^{12}$ the positivity itself needs no new computation: an off line pair makes the density negative only within $\frac12$ of its ordinate, so $\rho_\varepsilon>0$ on $[0,3\times10^{12}-\frac12]$ for every $\varepsilon>0$ (Corollary~\ref{cor:dbeta:window}); there the certificates test the evaluators against the verification and measure the margins.

Below height $200$ the certificate is within reach now. The ball arithmetic of Section~\ref{ch:polya} certifies the rounding only. Adding the discretisation bound of Section~\ref{sec:ch4:trapezoid} and of Proposition~\ref{prop:tilt}, bounds for the truncations of the theta series and of the $u$ axis, a bound on the change of the argument between grid points, and an explicit bound on the remainder of Proposition~\ref{prop:mass}, which is $O(\varepsilon^3)$ (Remark~\ref{rem:ch7:cubic}), turns positivity of $\rho_\varepsilon$ on $0<z\le200^2$ at each computed $\varepsilon$, the integer identity of Proposition~\ref{prop:zerodensity} and the HCM derivative signs of Section~\ref{sec:weak:hcm} into theorems on that window.

Above that window the evaluator changes. The tilted contour gives $25$ digits at height $5000$ and resolves the Lehmer pair near height $7005$ (Sections~\ref{sec:newcomp} and \ref{sec:ch5:lehmer}), but its number of nodes grows linearly with the height. Chirp transform batching of the P\'olya route, and the hybrid of Section~\ref{sec:ch5:hybrid}, a Riemann Siegel or Odlyzko Sch\"onhage evaluator at $\tfrac12+\varepsilon+i\theta$ with Odlyzko's zeros as the RH reference, would give the positivity margin, the matched total variation and the zero density scan near the $10^{12}$th \cite{Odlyzko87}, $10^{21}$st and $10^{22}$nd zeros \cite{Odlyzko01,OdlyzkoTables}. The same zeros would measure the cell constant of Section~\ref{ch:spacing} at large height, whose behaviour is open (Section~\ref{sec:ch15:spacings}), and allow a comparison of the cell errors with the prediction of the sine kernel process of the unitary ensemble.

\subsection{Calibration against known configurations}\label{sec:open:calib}

The diagnostics of Chapters~\ref{chap:measure} to \ref{chap:heat} are calibrated when they are run on configurations whose zeros are known exactly, and three such computations remain. Section~\ref{ch:deconv} ran the posterior predictive check of \cite[Section 19.3.7]{P26}, with a positivity prior on the measure $U$ in $\rho_\varepsilon=P_\varepsilon*U$, at $\varepsilon=0.1$ on one window, and calibrated it with a displaced zero; it is the linear inverse problem of \cite[Section 19.3]{P26}, whose data the P\'olya route supplies. Three extensions are open: smaller resolutions, where the signature of a displaced zero is largest and the P\'olya route supplies $\rho_\varepsilon$ to about $60$ digits below height $200$ (Section~\ref{sec:ch4:trapezoid}), whereas the deconvolution used double precision; continuous atom positions, which would remove the baseline of the grid; and larger heights. The square root law \eqref{eq:ch8:threshold} of Section~\ref{sec:ch8:calibration}, which follows from Proposition~\ref{prop:signature}, sets what the check can detect. The Williams benchmark of Section~\ref{sec:ch19:williams} needs no separate computation, since its distances follow from the closed forms of Section~\ref{sec:ch3:bench}.

The de Bruijn Newman flow of Section~\ref{ch:dbn} manufactures off line pairs in a controlled way. Running these Thorin diagnostics along the flow through the collision of the Lehmer pair near $7005$ would calibrate the zero density scan, the matched total variation and the negative mass against configurations that are known exactly, extending the collision time computed in Section~\ref{ch:dbn} to the full set of diagnostics.

Generalised prime systems of Beurling whose zeta functions provably have only real ordinates give positive Thorin measures by construction. Measuring how fast their clock transforms approach $1/G$, in the distances of Section~\ref{ch:distances} and along a sequence with an accumulation point as the closure criterion of Corollary~\ref{cor:weak:closure} \cite[Propositions 18.10 and 18.14]{P26} allows, is the computational side of that criterion; the Kent approximants of Section~\ref{sec:ch18:kent} are the first instance.

Two computations of Chapter~\ref{chap:probability} remain at their edges. The Brownian boundary of Section~\ref{ch:boundaries} reproduces the law of the clock in the bulk, to $0.0020$ in Kolmogorov distance for the boundary of the Volterra equation and to $0.0059$ for the tangent approximation of Roberts and Shortland (Section~\ref{sec:ch19:zones}), but the zone of large times $t$ is checked by no simulation, since only about one path in ten million survives past $t=0.1$. Importance sampling with an exponentially tilted drift, or multilevel splitting, would check that zone and the logarithmic correction of the asymptote of the boundary directly. The Cauchy integral method of Section~\ref{ch:li} extends to $n$ in the thousands, at a cost of about $0.03n$ digits lost to cancellation, so that the precision grows linearly in $n$. Comparing the P\'olya values of the Keiper Li coefficients with the sums $\sum_{\gamma>0}4\sin^2(n\arctan\frac1{2\gamma})$ over the zeros to large $n$ would test how the Li functional sees the zeros near height $n$.

\subsection{Testing stronger conjectures before pursuing them}\label{st:sec:wp5}

Conrey and Li showed that a positivity condition in de Branges' approach to RH is false even though RH may be true \cite{CL00}. Any intermediate statement proposed on the Thorin path should therefore be tested numerically first. The candidates are a lower bound $\rho_{\alpha-1/2}(\theta)\ge c(\alpha)\log\theta$ outside a sparse set, the first of the four questions of Section~\ref{ag:sec:exactly}; monotonicity or convexity of the margin in $\alpha$; positivity of the Thorin density for $\alpha$ slightly below one, uniformly on large windows; hyperbolic complete monotonicity of the Laplace transform $\varphi_\alpha$ of $T_\alpha$ below one; and variogram inequalities at single scales. The first two are partly settled: the lower bound holds for $\rho_{\alpha-1/2}$ itself outside a set of heights of density zero, although not near the zeros where Lemma~\ref{ag:lem:cancel} would need it (Remark~\ref{rem:dbeta:average}), and the margin divided by $\alpha-\frac12$ decreases in $\alpha$, with the ratio $\rho_{\alpha-1/2}(\theta)/(\alpha-\frac12)$ nonincreasing in $\alpha$ at every height outside the windows of the zeros off the line (Theorems~\ref{thm:dbeta:comp} and \ref{thm:dbeta:minall}). Each can be evaluated on the computed windows and at the heights of Odlyzko's tables, and any statement that fails there should be discarded. In parallel, the de Bruijn Newman flow, along which the reciprocal of the heat evolved $\Xi$ is a GGC for $t\ge0.2$ and RH is the statement at $t=0$ \cite{Polymath,PT21,RodgersTao}, gives a one parameter test bed in which the same diagnostics can be followed as $t$ decreases.

The variogram inequalities can be tested in finite dimension. The reformulation of Section~\ref{ch:sin2more} poses finite dimensional feasibility problems: is the regularised pole mass dominated by the prime and archimedean mass on the cone of variograms of positive definite functions of a given dimension? Solving them with explicit dual certificates at increasing dimension would measure how the margin behaves.

The cross checks of the fourth aim test the zeros themselves. Riesz's function, computed from the M\"obius function alone, sees only the first few zeros, and its twists by $(x/n^2)^{ih}$ move its sensitivity to the zeros near height $2h$ while keeping the criterion equivalent to RH (Proposition~\ref{ri:prop:twisted}); agreement of these twisted functions with the explicit formula over the computed zeros is a check that shares no code or method with the computation of the zeros (Remark~\ref{ri:rem:indep}). The open problems on the reciprocal zeta function (Section~\ref{rz:sec:open}), among them Conjecture S on the symmetrised truncations of the Biane, Pitman and Yor lattice and the Nyman Beurling problem, and the four questions of Section~\ref{ag:sec:exactly} are the natural targets for all these tests.

\subsection{Open problems from free probability}\label{sec:wigfree:open}

The note \cite{P26wig} closes with eight problems, and five of them are settled or taken up in the book. The positivity of the Thorin measure of $\Upsilon(T)$ and the infinite divisibility of $T$, classical and free, follow from Theorem~\ref{thm:ch12:idsd} through Theorem~\ref{thm:wigfree:clock}(a) and (b). The Gamma(2) problem is settled by Propositions~\ref{prop:wigfree:wishart} and~\ref{prop:wigfree:peel} and Theorem~\ref{thm:wigfree:nomix}. The failure curve of the Hankel test along the flow and the collision brackets are the subjects of Sections~\ref{sec:ch21:hankel} and~\ref{sec:wigbiane:brackets}. What remains open is the following.

\begin{problem}\label{prob:wigfree:jensen}
Jensen polynomials as finite free products. The multiplicative convolution $\boxtimes_d$ multiplies normalised coefficients and keeps roots real and of one sign (Lemma~\ref{lem:wigfree:laguerre}(i)), and the normalised coefficients of the reversed Jensen polynomial $x^dJ_{d,n}(1/x)$ are, up to the signs $(-1)^i$, the numbers $(n+i)!\,b_{n+i}$, $0\le i\le d$. A factorisation $(n+i)!\,b_{n+i}=\alpha_i\beta_i$ into the normalised coefficients of two polynomials with only real roots of one sign would therefore give real zeros of $J_{d,n}$. Find such factorisations for every shift $n$ at a fixed degree $d$, or show that none exists with factors of a given kind. The obvious split, a Gaussian multiplier sequence $q^{i^2}$, $0<q\le1$, times a residual close to the coefficients of $(1+x)^d$, fails, because a root of multiplicity $d$ splits into nonreal roots under a generic perturbation; and by Proposition~\ref{prop:wigfree:peel} the factor $1/(a)_i$ can be multiplied in but not divided out. A factorisation has to put real structure into both factors \cite{P26wig}.
\end{problem}

\begin{problem}\label{prob:wigfree:laguerre}
Laguerre and Meixner Pollaczek expansions. Theorem~\ref{thm:wigfree:nomix} excludes scale mixtures of gamma laws for $|X|$ and $X^2$, and Proposition~\ref{prop:wigfree:peel} shows that removing a gamma layer keeps zeros real while adding one does not. Find a variable in which the P\'olya density, or $\Xi$ itself, has an expansion in Laguerre or Meixner Pollaczek polynomials with positive coefficients and truncations with only real zeros \cite{P26wig}. The symmetric Meixner Pollaczek expansions of Section~\ref{rz:sec:mpexp} converge slowly, and a real rooted construction from their coefficients is a problem of Section~\ref{rz:sec:open}.
\end{problem}

\begin{problem}\label{prob:wigfree:jacobi}
Jacobi parameters. Under RH, where $\nu$ is a positive measure and all its Jacobi parameters exist, prove that $\alpha_n\sim2\beta_n\sim\frac12L_n$, with $L_n$ given by \eqref{eq:wigfree:abel} for the smooth count: a spectral form of the Riemann von Mangoldt formula. Table~\ref{tab:wigfree:jacobi} supports it to a few per cent for $n\le49$, but the local model of Section~\ref{sec:wigfree:jacobi} is a heuristic, outside the scaling of \cite{KVA99}. Decide between the rules $n(L_n)=n$ and $n(L_n')=n-\frac12$, which the data do not separate, and explain the scatter of the ratios, about six per cent, through the individual gaps between the zeros, as \cite{P26wig} proposes.
\end{problem}

\begin{problem}\label{prob:wigfree:window}
The order at which a collision is seen. Section~\ref{sec:ch21:hankel} computes the failure times $t^*_N$ of the Hankel test along the flow to $N=50$ and explains their steps by collisions; the rule of \cite{P26wig} that the collision of $(\gamma_k,\gamma_{k+1})$ is registered from order $2k+2$ on fails at $N=50$ (Table~\ref{tab:ch21:nfail}). Find the rule that gives the order at which a pair is registered, as a function of the index $k$ of the pair and of its distance from the axis, and follow the staircase of $t^*_N$ to orders of a few hundred, towards the limit $0$ that it has under RH.
\end{problem}

\begin{problem}\label{prob:wigfree:collide}
Unconditional collision brackets. Section~\ref{sec:wigbiane:brackets} brackets the collision times by quantities along the flow (Corollaries~\ref{cor:wigbiane:delta} and~\ref{cor:wigbiane:kappa}) and through the nonreal zeros (Theorems~\ref{thm:wigbiane:local} and~\ref{thm:wigbiane:upper}), and checks the hypotheses on a grid of times for the $58$ collisions of the census. Find the a priori control of the nonreal zeros of $H_t$, $t<0$, asked for in Problem~\ref{prob:wigbiane:control}, which would make the brackets statements about $H_t$ computed from $H_0$ alone.
\end{problem}

\begin{problem}\label{prob:wigfree:ladder}
Moving the ladder without new zeros. Theorems~\ref{thm:ch12:cmint} and~\ref{thm:ch12:window} prove the derivative signs $(-1)^kW^{(k)}>0$ on ranges that grow with the height to which RH is verified (Section~\ref{sec:ch12:reach}), and complete monotonicity of $W$ on $(0,\infty)$ is RH. Find a mechanism that extends these ranges with no new information on the zeros, for instance a functional equation for the law of $T$ under $\Upsilon$, as \cite{P26wig} suggests. Theorem~\ref{thm:wigfree:clock}(c) restricts the form such an equation can take: $T$ is not the Upsilon image of a law in BO, and under RH it is the image of the law $\mu_0$, which is neither a GGC nor in BO.
\end{problem}

\subsection{Lessons from the computations}\label{ch:lessons}

The computations of the book distil into lessons for anyone working on the Thorin route to RH. Each is a statement that a computation in the book established or refuted, and each saves effort on a route that cannot succeed or points to one that might.

Equivalent statements are not equally computable. The Thorin condition has at least eight equivalent faces: positivity of the Thorin measure, complete monotonicity of the heat trace, Stieltjes form of $G'/G$, Hankel positivity at one time, HCM of the clock, Lagarias positivity, Weil positivity and Li positivity. Mathematically they are the same statement. Computationally they differ by many orders of magnitude in conditioning. The heat trace at moderate time loses about $87t$ digits to cancellation (Section~\ref{ch:arith}); the normalised minors of the even Hankel block at $t=0.05$ fall to $10^{-44}$ by size five (Section~\ref{ch:heat}); the Keiper Li coefficients weight a zero at height $\gamma$ by $n^2/\gamma^2$ until $n$ reaches $\gamma$ (Section~\ref{ch:li}); the Thorin density at resolution $\varepsilon$ is local in the ordinate and sees each zero at full size $1/\varepsilon$ (Section~\ref{ch:distances}). A computational attack should choose its face by conditioning, and the local face is the Thorin density.

The primes compute the wrong end. The sine squared lemma writes the Thorin density as a sum of positive prime and archimedean pieces and one negative pole piece. It is the right way to understand the structure: the pole is the only term of the wrong sign, and RH is the statement that it is dominated after the kernel acts; Section~\ref{sec:ch16:sin2} shows what changes for an $L$ function without a pole. It is the wrong way to compute near the line, where $k$ digits at resolution $\varepsilon$ need primes up to about $10^{2k/\varepsilon}$ even when no zero lies to the right of the line (Section~\ref{sec:ch10:selberg}). Any argument that proposes to establish positivity near the line by bounding prime sums must overcome that exponential, and Section~\ref{ch:rz3} shows the same obstruction in another form: the bias that keeps the canonical function positive is of order $x^{-1/2}$, while pointwise bounds on the primes control it only to order $x$ (Proposition~\ref{rz:prop:bias}).

The heat trace meets the same obstruction in its own terms. By Proposition~\ref{prop:ch12:gauss}, RH is the positivity of the Gaussian smoothing $\mathcal G(u,t)$ of the zeros at every height $u$ and every $t>0$, and \eqref{eq:ch12:Garith} writes the symmetric smoothing $\mathcal G(u,t)+\mathcal G(-u,t)$ as a polar term, negligible at large $u$, plus an archimedean term $A(u,t)$, close to $\log(u/2\pi)/(2\sqrt{2\pi t})$, minus a prime sum $P(u,t)$. Bounding $|\cos(u\log n)|$ by one bounds $P(u,t)$ by $S(t)/\sqrt{2\pi t}$, where $S(t)=\sum_{n\ge2}\Lambda(n)n^{-1/2}e^{-(\log n)^2/8t}$ is, up to a factor, the prime term of \eqref{eq:ch2:Warith} at time $2t$, and so proves $\mathcal G(u,t)+\mathcal G(-u,t)>0$ for $t$ below the time $t_0(u)$ at which the bound meets the archimedean and polar terms, which agrees to five digits with the root of $2S(t)=\log(u/2\pi)$ (computed by \texttt{phasefree.py}): $t_0$ is $0.45$ at height $10^3$, $1.72$ at the verified height $3\times10^{12}$ and $3.64$ at $10^{50}$, to the digits shown (proved in ball arithmetic by \texttt{phasefree\_arb.py}), and it grows only doubly logarithmically in $u$, since $S(t)$ is asymptotic to $\sqrt{8\pi t}\,e^{t/2}$. For $\mathcal G(u,t)$ itself the bound must also exceed $|\mathcal G(-u,t)|$, at most $e^{t/2}((u+14)^2+\frac1{2t})e^{-2t(u+14)^2}$ by Section~\ref{sec:ch12:ray}, which is below $10^{-880}$ for $u\ge10^3$ and $t\ge10^{-3}$; on that range the time at which positivity of $\mathcal G(u,t)$ is proved differs from $t_0(u)$ only far beyond the digits shown. These resolutions are too coarse to see the quadruples of Table~\ref{tab:ch13:first}: at their heights $t_0$ lies between $0.311$ and $0.511$, while within $8$ of each moved pair they make $\mathcal G$ negative only for $t>1.18$ (both proved in ball arithmetic). Made uniform in the order, the same bound proves the signs of all derivatives of order above $10^{17}$ for $t\le1.29$, which with the verified zeros makes $W$ completely monotone there (Theorem~\ref{thm:ch12:cmint}). That passes $1.18$ only because, along the rays, those orders lie at heights above $2.78\times10^8$, where $t_0$ exceeds $1.29$; at the heights of the quadruples the theorem uses the verified zeros instead. Nor can any bound on $P$ that holds uniformly in the height do better, since by Dirichlet's theorem on simultaneous approximation the phases $u\log n$ of any finite set of prime powers come arbitrarily close to multiples of $2\pi$ at arbitrarily large $u$, where $P(u,t)$ therefore comes arbitrarily close to $S(t)/\sqrt{2\pi t}$. Nor does positivity at a single resolution give a zero free strip: at fixed $t$ the bound $S(t)/\sqrt{2\pi t}$ on $P$ proves $\mathcal G(u,t)+\mathcal G(-u,t)>0$ at every height $u$ with $t_0(u)>t$, that is above about $2\pi e^{2S(t)}$, which is $8.1\times10^{212}$ at $t=6$, and $\mathcal G(u,6)>0$ itself for every $u\ge10^{213}$ (proved in ball arithmetic, with the bound on $|\mathcal G(-u,6)|$ above), whatever the zeros are, so a strip needs $t$ to grow with the height. A proof along the heat trace must use the phases, uniformly in the height.

Limits do not create positivity. Weak limits of Thorin measures preserve positivity but cannot manufacture it. The Cauchy flow converges vaguely to the measure of the line zeros whether or not RH holds, because an off line pair at distance $d$ from the line contributes to $\mu_\varepsilon$ a signed measure whose negative part, for $\varepsilon<d$, has mass only $\frac2\pi\arcsin\frac\varepsilon d\approx2\varepsilon/(\pi d)$, that is $\varepsilon/(\pi d)$ per displaced zero (Proposition~\ref{prop:signature}(iii) and Section~\ref{ch:weak}). Positivity must therefore be proved at every resolution, not in the limit. The closure criterion helps only for approximants whose positivity is proved independently, such as the Beurling systems and the Kent approximants of the first mechanism in Section~\ref{sec:ch25:open}. Verification removes the factors of the verified zeros from the question and leaves it for the tail factor of the clock \cite[Corollary 18.11]{P26}. A proposal that reaches the Thorin condition as a limit must say where the positivity of its approximants comes from.

\subsection{Closed routes and the reach of finite information}

Multipliers and reformulations do not help by themselves, and several routes are now closed or reduced to RH itself. Multiplying the clock by any GGC leaves the problem exactly as hard (Proposition~\ref{rz:prop:nomult}). HCM of the reciprocal Xi function is equivalent to the GGC property by Bondesson's theorem, and so to RH (Corollary~\ref{cor:hcm}); it is not a way round it. Weak closure of the Cauchy flow cannot extend the family (Section~\ref{sec:weak:limits}). The clock is not a Brownian first passage time to a random level for any constant drift, so the simplest hitting time representations are excluded unconditionally (Proposition~\ref{prop:nogo}), although at every basepoint above the centre the shifted clock is a subordinated hitting time (Section~\ref{ch:bphit}). The Bernstein Pick route through the moments of the P\'olya density is closed by a single interval certificate (Section~\ref{ch:finite}). And the prime side of the sine squared lemma cannot compute near the line (Section~\ref{ch:sin2more}). Each of these negative results was obtained cheaply, and each removes a route that had looked promising. Sections~\ref{ch:riesz} and \ref{ch:sato} close two more: the exact representations of the Riesz coefficients through the clocks cannot supply Riesz's growth bound (Section~\ref{ri:sec:help}), and the theorems of Sato and Bondesson reach nothing below the basepoint one (Section~\ref{sb:sec:cannot}). Section~\ref{sec:pakes} closes three more: the clock is not the martingale limit of a supercritical continuous state branching process (Theorem~\ref{pk:thm:nobranch}), its index among Urbanik's classes is decided with no hypothesis (Theorem~\ref{pk:thm:ladder}), and no property that the damped kernels $e^{tv^2/4}\Phi(v)$, $t<0$, share with $\Phi$ can imply RH (Proposition~\ref{pk:prop:filter}).

Finite information has a precise reach. Every finite computation in the book establishes a statement on a window and nothing beyond it. The perturbation theorem makes this exact for the heat trace tests: for every $N$ and every sufficiently large height $R$ for which RH has been verified up to $R+50$, there are functions with off line zeros just above $R$ that keep the first $N$ Hankel blocks positive and the derivative signs through order $10^{17}$ (Theorem~\ref{thm:ch13:perturb} and \cite[Theorem 11.6 and Proposition 11.8]{P26}), and every such result in the book is of that kind. The local visibility condition makes it exact for positivity at a fixed small resolution: to first order in $\varepsilon$, an off line zero at distance $d$ from the line makes $\rho_\varepsilon$ negative at its ordinate if and only if $2/d^2>\sum_j\delta_j^{-2}$, where $\delta_j$ are the distances to the other ordinates, and the condition does not improve as $\varepsilon\downarrow0$ (Section~\ref{sec:ch13:visibility}). For neighbours at the mean spacing near height $T$ it fails once $d$ exceeds about $0.78$ of the spacing, which a zero can do only above height about $1.1\times10^5$ (Section~\ref{ch:spacing}); such a zero shows itself only at resolutions comparable to $d$. The zero density identity of Proposition~\ref{prop:zerodensity} is the most complete finite test, because it scans every resolution at once and returns integers, but it too is a statement at a finite height. What finite computation buys is a calibration of how much of RH a given amount of work establishes, and that calibration is itself useful: it tells a proof what it must supply.

The perturbation theorem also gives a test that a proposed proof of RH must pass. Replace $\xi$ by the function $F$ of Theorem~\ref{thm:ch13:perturb}, with $R$ above every height that the proof uses explicitly. A step that holds for $F$ as well as for $\xi$ cannot be the step that excludes zeros off the line, because $F$ has four. The steps that survive include the Riemann von Mangoldt formula with an explicit error, which $F$ satisfies with the error enlarged by one on $[A,B)$, the only heights where the two counts differ; the verified zeros below $R$; the small time expansion of the heat trace through the constant term, since $W_F-W=\big(\frac12(B-A)^2+\frac18\big)t+O(t^2)$; infinite divisibility, self decomposability and log convexity; the Hankel blocks of size $N$ when $R\ge T_N$ and the derivative signs up to order $10^{17}$ when $R\ge10^{12}$; and the order, type and functional equation. What separates $\xi$ from $F$ is arithmetic, the Euler product and the explicit formula with the prime weights $\Lambda(n)$, together with local counts of the zeros on the line in windows that contain the height $g$ of the quadruple (Section~\ref{sec:ch13:perturb}). A proof through the derivative signs can do without them only below the order at which the signs of $F$ first fail, which is about $t^*R^2$, with $t^*$ between $1.18$ and $4.03$ in the cases computed below height $2000$ and, by the model of Section~\ref{sec:ch13:perturb}, of order one or more at larger heights, far beyond $10^{17}$. The function $F$ exists unconditionally when $R+50\le3\times10^{12}$; for larger $R$ the test still applies, because a correct proof of RH makes the hypothesis of Theorem~\ref{thm:ch13:perturb}(b) true for every $R$. The criticality filter of Section~\ref{pk:sec:filter} is the companion test for properties of the P\'olya kernel: a step that holds for $e^{tv^2/4}\Phi(v)$ with some $t<0$, whose transform has nonreal zeros by the theorem of Rodgers and Tao, cannot be the step that excludes zeros off the line.

\subsection{What a proof must supply}

What remains is a single positivity statement, in equivalent forms: the Thorin density $\rho_\varepsilon$ is nonnegative at every basepoint above the centre; the Gaussian smoothing $\mathcal G(u,t)$ of the zeros is nonnegative at every height and every resolution (Proposition~\ref{prop:ch12:gauss}); the shifted clock is a GGC at every basepoint above the centre, which by Theorem~\ref{ri:thm:ggcriesz} is also Riesz's criterion; the pole mass is dominated by the prime and archimedean mass on every variogram of a positive definite function; the zero density $\nu(\tfrac12+\varepsilon,T)$ vanishes for every $\varepsilon$ and $T$. A proof of the Thorin condition must supply this positivity uniformly in the height and in the resolution, and it needs a structural reason for it. The candidates consistent with everything computed here are the three mechanisms for step~8 in Section~\ref{sec:ch25:open}: a positive approximating family with independently proved positivity converging under the closure criterion, an operator whose spectrum is the set of ordinates, and a direct inequality for the variogram pairing. The computations of this book cannot supply any of these. What they can do, and have done, is test candidate lemmas quickly, produce counterexamples, exclude routes that fail, and measure the margins that a successful argument will have to respect.

The programme of this section has one organising principle: compute the quantities whose positivity or growth is equivalent to RH, at the places where a proof would have to control them, and compare what is found with what random matrix theory and the explicit formula predict. Riesz's criterion in its classical form is blind above the first few zeros; twisted Riesz functions repair this and give an arithmetic filter bank independent of evaluations of $\zeta$ (Section~\ref{ri:sec:filters}). Thorin densities at small resolution can be certified on windows by ball arithmetic and the zero density identity (Section~\ref{st:sec:wp2}). Below height $2198$ their positivity margins agree with the circular unitary ensemble in the median to $1\%$ and lie slightly above it in the lower tail, which makes the room RH has at each height predictable, and the explicit defect of the almost GGC factorisation measures how far each basepoint is from its target (Sections~\ref{ag:sec:margins} and \ref{ag:sec:factor}). None of this proves RH; all of it sharpens the target and tests the statements a proof would use. Appendix~\ref{app:code} lists the scripts behind the computations of the book.
